% --- Back cover ---
% Large PERSONIX wordmark, no logo image. Cream paper inherited globally.
% Absolute (tikz overlay) positioning so the footer sticks to the bottom.
\clearpage
\thispagestyle{empty}
\begin{tikzpicture}[remember picture, overlay]
  % --- Wordmark ---
  \node[anchor=north, font=\sffamily\bfseries\fontsize{48}{52}\selectfont]
    at ([yshift=-26mm]current page.north) {PERSONIX};
  \node[anchor=north, font=\rmfamily\itshape\large]
    at ([yshift=-44mm]current page.north) {Безкомпромисна промяна};
  \node[anchor=north, font=\rmfamily\small,
        text width=\paperwidth-50mm, align=center]
    at ([yshift=-54mm]current page.north)
    {Нецензурируема и неподкупна децентрализирана репутационна мрежа};
  % --- Body ---
  \node[anchor=north, text width=\paperwidth-50mm, align=justify, font=\small]
    at ([yshift=-72mm]current page.north)
    {%
      Държавата работеше, когато комуникацията беше бавна, решенията бяха
      локални, а властта живееше в същия град, който управляваше. Днешните
      мрежи обръщат и трите допускания --- и институциите, изградени върху тях,
      вече не пасват на света, който управляват. Тази книга описва инструмент,
      който пасва: децентрализирана репутационна мрежа, в която идентичността
      е твоя, за да я запазиш, истината е публично одитируема, а подкупът е
      репутационно самоубийство.

      \smallskip
      Не манифест. Не прогноза. Работещ план как може да бъде изграден
      наследникът на държавата --- доброволно, една декларация след друга.
    };
  % --- Footer (anchored to the bottom of the page) ---
  \node[anchor=south, font=\sffamily\footnotesize,
        text width=\paperwidth-50mm, align=center]
    at ([yshift=12mm]current page.south)
    {%
      Pavel Kudrna\quad\textbullet\quad Прага, 2026\par
      \href{https://personix.org}{personix.org}\par
      Лицензирано под Creative Commons Attribution-ShareAlike 4.0 International
      (CC BY-SA 4.0).
    };
\end{tikzpicture}%
\mbox{}%
\clearpage
